在“双碳”目标和配送集约化要求下，多车场混合车队既要通过跨场协同降低路径成本，又需依据时变电网碳强度安排电动车充电，并保证各车场愿意参与；动态订单进一步增加了方案执行难度。针对上述问题，建立可行配送趟—实体车排班两层路径优化模型，以运营成本和碳结算成本之和最小为目标，刻画多趟衔接、连续充电、成员参与和滚动状态继承。设计含跨场重组算子的自适应大邻域搜索，并在固定排班内进行碳感知充电重调度。通过区域配送网络、连续电网日和动态事件流检验模型与算法。结果表明，充电择时是路径与车型减排的补充，协同价值受客户空间组织和参与条件共同约束；滚动重规划还需同时满足状态可继承性和计算时限。
